\chapter{Gout Medication}
\index{Gout Medication}

\section*{Definition and Overview}
Gout medication falls into two groups: medicines that treat acute attacks and medicines that lower uric acid levels to prevent future attacks. Acute-attack medicines, including non-steroidal anti-inflammatory drugs (NSAIDs), colchicine, and corticosteroids, rapidly reduce pain and inflammation and are most effective when started early. Urate-lowering therapy, principally allopurinol, reduces the body's uric acid levels over weeks to months, so that crystals dissolve and attacks stop. Urate-lowering therapy is taken long term, usually for life, and is started after the acute attack has settled. Used correctly, gout medication prevents attacks, stops joint damage, and allows people to live free of gout.

\section*{Epidemiology}
Gout affects around 4--6\% of adults, and medication is central to its management. Acute attacks are treated in most cases with NSAIDs or colchicine, and around a third of people with gout are prescribed urate-lowering therapy. However, many people who would benefit from urate-lowering therapy are not on it, and adherence is often poor, with studies suggesting only around half take it as prescribed. Suboptimal treatment leaves many people with recurrent attacks and preventable joint damage, which is why education about gout medicines is so important.

\section*{Aetiology and Pathophysiology}
The medicines used for gout target the two processes driving the disease.

\begin{itemize}
    \item \textbf{Inflammation:} acute attacks are caused by urate crystals triggering an intense inflammatory response; NSAIDs, colchicine, and corticosteroids dampen this response
    \item \textbf{Uric acid production:} the body produces uric acid as it breaks down purines; allopurinol and febuxostat block the enzyme xanthine oxidase that produces it
    \item \textbf{Uric acid excretion:} most people with gout under-excrete uric acid; uricosuric medicines such as probenecid increase excretion but are used less often
    \item \textbf{Crystal dissolution:} when uric acid levels fall below the saturation point, existing crystals slowly dissolve
    \item \textbf{Flare on starting treatment:} lowering uric acid quickly can mobilise crystals and trigger an attack, which is why low-dose colchicine or an NSAID is used during the first months
    \item \textbf{Target levels:} treatment aims for a serum uric acid below about 0.36 mmol/L, and lower (0.30) if tophi are present
\end{itemize}

\section*{Clinical Features}
Gout medicines are used in specific clinical situations.

\begin{itemize}
    \item An acute attack: sudden severe joint pain, redness, and swelling needing rapid treatment
    \item Recurrent attacks: two or more attacks a year indicating the need for urate-lowering therapy
    \item Tophi, chronic gouty arthritis, or kidney stones related to uric acid
    \item High uric acid levels in people with gout, even between attacks
    \item Prevention of flares during the first months of urate-lowering therapy
    \item People with coexisting conditions, such as kidney disease or peptic ulcers, who need tailored medicine choices
\end{itemize}

\section*{Differential Diagnosis}
Medicine choices are individualised, and other factors are considered before prescribing.

\begin{itemize}
    \item Kidney function, which affects the choice and dose of urate-lowering therapy
    \item Gastrointestinal disease, including ulcers and inflammatory bowel disease, which affects NSAID and colchicine use
    \item Cardiovascular disease, which is associated with gout and affects NSAID choices
    \item Drug interactions with warfarin, diuretics, and other medicines
    \item Pregnancy and breastfeeding, which limit some options
    \item Allergy to allopurinol, which occurs in a small number of people and can be severe
    \item Distinguishing gout from septic arthritis, which needs urgent different treatment
\end{itemize}

\section*{When to Seek Medical Care}
Anyone with gout should be reviewed by a doctor to plan acute and long-term treatment. Seek medical advice if attacks become frequent, if medicines cause side effects, or before starting any new medicine, because some interact with gout treatments.

\textbf{Contact your local emergency services for:}
\begin{itemize}
    \item A hot, swollen joint with fever, which may be septic arthritis
    \item A severe rash, especially with mouth ulcers, fever, or facial swelling, within weeks of starting allopurinol, which may indicate a serious allergic reaction (DRESS or Stevens-Johnson syndrome)
    \item Severe vomiting, diarrhoea, or abdominal pain after colchicine, which can be toxic
    \item Signs of internal bleeding, such as black stools or vomiting blood, while taking an NSAID
\end{itemize}

\section*{Diagnosis and Investigations}
Medication is guided by diagnosis and monitoring.

\begin{itemize}
    \item \textbf{Confirmation of gout:} clinical features, joint fluid analysis, and blood uric acid levels
    \item \textbf{Baseline tests:} kidney function and liver function before starting urate-lowering therapy
    \item \textbf{Uric acid monitoring:} blood levels are checked regularly and the dose adjusted to reach the target
    \item \textbf{Full blood count:} occasional monitoring is needed with some medicines
    \item \textbf{Review of coexisting conditions and medicines:} ensures safe prescribing
    \item \textbf{HLA-B*5801 testing:} considered before allopurinol in people of Han Chinese, Thai, and Korean descent, who are at higher risk of severe allergic reactions
\end{itemize}

\section*{Management}
Gout medication is managed in a stepwise, planned way.

\begin{itemize}
    \item \textbf{Treat the acute attack early:} start an NSAID, colchicine, or corticosteroid within hours of symptoms for best effect
    \item \textbf{Start urate-lowering therapy once the attack settles:} usually allopurinol, beginning at a low dose and increasing gradually
    \item \textbf{Protect against flares:} take low-dose colchicine or an NSAID for the first three to six months of urate-lowering therapy
    \item \textbf{Treat to target:} adjust the dose until uric acid is below the target level
    \item \textbf{Continue long term:} urate-lowering therapy is lifelong for most people, and stopping it causes attacks to return
    \item \textbf{Manage coexisting conditions:} treat high blood pressure, kidney disease, and obesity alongside gout
    \item \textbf{Specialist referral:} for complex gout, tophi, or treatment failure
\end{itemize}

\section*{Complications}
\begin{itemize}
    \item Side effects of NSAIDs, including stomach ulcers, kidney injury, and increased cardiovascular risk
    \item Side effects of colchicine, including diarrhoea, nausea, and, at high doses, serious toxicity
    \item Side effects of corticosteroids, including weight gain, diabetes, and infection risk
    \item Allergic reactions to allopurinol, which can be severe and life-threatening
    \item Flares when urate-lowering therapy is started without flare prophylaxis
    \item Poor adherence, leading to recurrent attacks and joint damage
    \item Drug interactions with common medicines
\end{itemize}

\section*{Prognosis and Outcomes}
Gout medication is highly effective when used correctly. Acute-attack medicines settle attacks within days, and urate-lowering therapy stops attacks within several months in most people. Long-term treatment dissolves tophi, prevents joint damage, and restores normal joint function. People who reach their uric acid target and stay on treatment have an excellent prognosis and can expect to live free of gout attacks. The main reasons for poor outcomes are undertreatment, stopping medicines, and delayed treatment of attacks.

\section*{Prevention and Self-Care}
\begin{itemize}
    \item Keep acute-attack medicines at home and start treatment at the first sign of an attack
    \item Take urate-lowering therapy daily, at the same time, and never run out
    \item Attend blood tests to keep uric acid in the target range
    \item Take flare-prevention medicine during the first months of urate-lowering therapy
    \item Report side effects promptly rather than stopping medicines on your own
    \item Tell your doctor about all other medicines you take
    \item Carry a list of your gout medicines and doses
    \item Combine medication with lifestyle measures: limit alcohol and purine-rich foods, maintain a healthy weight, and stay hydrated
\end{itemize}

\section*{Sources and Further Reading}
The following reputable sources provide further information for healthcare students and patients seeking reliable, current advice about gout medication.

\begin{itemize}
    \item Australian Rheumatology Association. \textit{Gout treatment information.} https://rheumatology.org.au/
    \item NPS MedicineWise. \textit{Medicines for gout.} https://www.nps.org.au/
    \item Healthdirect Australia. \textit{Gout: treatment with medicines.} https://www.healthdirect.gov.au/
    \item Arthritis Australia. \textit{Gout and its treatment.} https://arthritisaustralia.com.au/
    \item NHS. \textit{Gout: treatment.} https://www.nhs.uk/
    \item American College of Rheumatology. \textit{Gout management guidelines.} https://rheumatology.org/
\end{itemize}
